% !TEX encoding = UTF-8
% !TEX program = pdflatex
% !TeX spellcheck = en_GB
% !BIB = biber


\documentclass[english]{article}
\usepackage{babel}
\usepackage[utf8]{inputenc}
\usepackage{graphicx}
\graphicspath{{./images/}}
\usepackage{hyperref}
\hypersetup{
    colorlinks=true, 
	linkcolor=blue, 
	filecolor=blue, 
	citecolor = black,       
	urlcolor=blue, 
}

\usepackage{biblatex}
\addbibresource{essay_biblio.bib}

\title{An Essay for the exam of Network Security}
\author{A. Student}

\begin{document}

\maketitle

\begin{abstract}
The essay should be about 15 pages. As a suggestion, pretend that I am your future CTO, and have asked you to study a given problem and present me some solution. Your aim is to brief me about that solution, explain pros and cons, in order to be able to take an informed decision.
\end{abstract}


\tableofcontents
\newpage

\section{Introduction}


In the introduction, give a description of the subject: applications, importance, aims, etc.

I suggest using Bib\TeX\ or Bib\LaTeX\ and \texttt{biber} for bibliography. Reference are easily cited like this: \cite{stallings2015computer}.

\section{Threat model}
Describe the potential threats, vulnerabilities, and risks to a computer network or system. 
%It involves assessing the various ways in which an attacker could compromise the system's confidentiality, integrity, or availability.
More precisely, define:
\begin{itemize}
\item Assets: the valuable resources, data, or systems which are to be protected.
\item Attackers: the potential threat actors %their motivations, capabilities, and the resources they may have at their disposal.
\item Attack vectors: Identify the potential ways in which attackers could exploit vulnerabilities to gain unauthorized access or cause harm. %This includes considering both external threats (such as network attacks) and internal threats (such as insider attacks).
\end{itemize}

\section{Security goals}
Describe the security goals that we aim to achieve on the identified assets, according to the CIAAA model.


\section{Security service and implementation}
Describe the security services to be implemented, and how these are implemented: protocols, algorithms, procedures, etc.


\section{Attacks and Vulnerabilities}
Possible vulnerabilities (if known) and corrections (if known)


\section{Conclusions}
Critical analysis about the effectiveness of the solution.

\appendix
\section{Something I forgot}
If needed, extra material (eg. long code listings, traces, logs, etc) can be put in the appendix.


\printbibliography
\end{document}
